% explainer-source-hash: 57bcbdec561bab334481f03dc7f1b4d5b8f071f1c30602f865e48dc1657f0e38
% explainer-source-commit: 40e6ab02e9e59f22535727fa3f73442500a09963
% explainer-generated: 2026-07-16
% Regenerate with /explain-all (see WIRING.md). Do not edit this stamp by hand:
% it is how the toolchain knows whether this document still matches the code.
\documentclass[11pt,a4paper]{article}
\input{preamble}

\title{\textbf{Hi-Tec Shopify Automation}\\[0.3em]\large A Deep Dive}
\author{Portfolio explainer, built from the repository source}
\date{}

\begin{document}
\maketitle
\thispagestyle{empty}

\begin{keyidea}{The whole project in one paragraph}
A bearings retailer sells thousands of products through an online shop hosted by
Shopify. The quantities and costs they actually trust live in a different
spreadsheet, exported from their warehouse system. Nobody had a way to push one
into the other, so a person did it by hand. This project is three small Python
programs that do that reconciliation on the exported spreadsheet files. They
never contact the shop over the network. They read a file, add columns, and
write a new file that a human reviews before uploading it back.
\end{keyidea}

\tableofcontents
\newpage

\section{Context and the constraint that shaped everything}

\subsection{Who this was for}

Hi-Tec Bearings is a retailer of bearings (the metal rings that let a wheel or a
shaft spin with low friction). They run a shop on Shopify. Their catalogue has
thousands of product variants. This was freelance work done solo by Salman Adnan
in 2023.

\begin{jargon}{Shopify}
A hosted e-commerce platform. A merchant does not run their own web server; they
pay Shopify, and Shopify runs the storefront, the checkout and the product
database. The merchant administers it through a web control panel called the
\textbf{Shopify admin}.
\end{jargon}

\begin{jargon}{Product variant}
One purchasable version of a product. A bearing product \code{22210} might exist
in a \code{22210 K} version (tapered bore) and a \code{22210 E} version
(different cage). Each is a variant. In a Shopify export, each variant is its own
row, and the rows of one product share a \code{Handle}.
\end{jargon}

\subsection{The problem}

Two spreadsheets described the same stock and disagreed:

\begin{itemize}
  \item The \textbf{stock list}, exported from the client's internal warehouse
  system. It carries the quantity actually on the shelf and the cost actually
  paid.
  \item The \textbf{Shopify export}, exported from the Shopify admin. It carries
  what the storefront currently tells customers.
\end{itemize}

Keeping the second in step with the first was manual and took hours per update
cycle. A third irritation: useful facts about each bearing (its bore diameter,
outside diameter and width) existed only as an HTML table buried inside the
product description text, so the storefront could not filter or sort on them.

\subsection{The constraint: no API}

\begin{keyidea}{Everything here is offline file processing}
Nothing in this project talks to Shopify over the network. There is no API
client, no access token, no shop domain, and \file{.env.example} declares
plainly that no environment variables are needed. The tools operate on files
that a human has already exported.
\end{keyidea}

\begin{jargon}{API}
Application Programming Interface. A machine-to-machine doorway into a service.
Shopify does publish one, and a program with an access token could read and write
products directly over the network. This project deliberately does not use it.
\end{jargon}

The working loop the client actually used:

\begin{figure}[H]
\centering
\resizebox{\textwidth}{!}{
\begin{tikzpicture}[node distance=8mm and 12mm]
  \node[extbox] (store) {Shopify store\\(hosted, remote)};
  \node[box, right=of store] (exp) {Human exports\\from Shopify admin};
  \node[store, right=of exp] (csv) {Shopify export\\CSV / XLSX};
  \node[box, right=of csv, fill=softgrey] (tool) {One of the three\\tools runs};
  \node[store, below=14mm of tool] (out) {\code{updated\_*} file\\(new columns added)};
  \node[dec, left=of out] (rev) {Human\\reviews};
  \node[box, left=of rev] (imp) {Human re-imports\\via Shopify admin};

  \node[store, above=9mm of tool] (stock) {Internal stock list\\CSV / XLSX};

  \draw[flow] (store) -- (exp);
  \draw[flow] (exp) -- (csv);
  \draw[flow] (csv) -- (tool);
  \draw[dflow] (stock) -- (tool);
  \draw[flow] (tool) -- (out);
  \draw[flow] (out) -- (rev);
  \draw[flow] (rev) -- node[lbl, above]{approved} (imp);
  \draw[flow] (imp) -| (store);
\end{tikzpicture}}
\caption{The round trip. The tools occupy one box. The human occupies four, and
one of them is a review gate that no catalogue change can skip.}
\end{figure}

\begin{keyidea}{The constraint bought something real}
Because the tools cannot write to the shop, a wrong result is a wrong spreadsheet,
not a wrong storefront. The human review step in the middle is not ceremony; it
is the only thing standing between a bad match and a live price. A design that
used the API would have had to add that safety back deliberately.
\end{keyidea}

The cost of the constraint is equally real: without the API there are no stable
product identifiers to join on. The tools can only join on the text fields the
export happens to contain. Section \ref{sec:matching} is entirely about the
consequences of that.

\section{Repository layout}

\begin{figure}[H]
\centering
\begin{forest} dirtree
[hitec-shopify-automation
  [README.md]
  [LICENSE {\rmfamily\scriptsize (PolyForm Strict 1.0.0)}]
  [requirements.txt {\rmfamily\scriptsize (pandas, openpyxl)}]
  [.env.example {\rmfamily\scriptsize (declares: nothing needed)}]
  [quantity\_pricing\_updater/
    [update\_quantity\_pricing.py {\rmfamily\scriptsize (69 lines, CLI)}]
    [update\_quantity\_pricing.spec {\rmfamily\scriptsize (PyInstaller)}]
  ]
  [dimensions\_extractor/
    [dimensions\_extractor.py {\rmfamily\scriptsize (83 lines, CLI)}]
    [dimensions\_extractor.spec {\rmfamily\scriptsize (PyInstaller)}]
  ]
  [stock\_updater\_gui/
    [shopify\_stock\_updater.py {\rmfamily\scriptsize (354 lines, tkinter, no .spec)}]
  ]
  [sample\_data/
    [sample\_stock\_list.csv {\rmfamily\scriptsize (8 stock rows)}]
    [sample\_shopify\_stock.csv {\rmfamily\scriptsize (6 product rows)}]
  ]
  [docs/
    [screenshots/app.png]
    [explainers/ {\rmfamily\scriptsize (this document)}]
  ]
]
\end{forest}
\caption{Every tracked file. Fifteen of them. This is a small repository and the
document is sized to match.}
\end{figure}

The shape is one folder per tool plus shared sample data. There is no package,
no \file{\_\_init\_\_.py}, no shared module, and no test directory. Each script
is standalone and run directly.

\section{The two data files}

Understanding the tools requires understanding the two spreadsheet shapes they
join. Both are given as small fabricated samples in \file{sample\_data/}.

\begin{jargon}{CSV and XLSX}
\textbf{CSV} (comma-separated values) is a plain-text table: one line per row,
fields separated by commas, quotes around fields that contain a comma or a
newline. \textbf{XLSX} is the modern Excel binary workbook format. Both hold a
table; CSV is readable in any text editor, XLSX is not.
\end{jargon}

\subsection{The stock list}

Seven columns. The ones the tools use are in bold.

\begin{center}
\begin{tabular}{@{}ll@{}}
\toprule
\textbf{Column} & \textbf{Role} \\
\midrule
\code{No} & Row number. Unused. \\
\textbf{\code{Item}} & \textbf{The bearing designation, e.g. \code{6205} or \code{22210 K}. Join key.} \\
\textbf{\code{Brand}} & \textbf{\code{HI-TEC} or a competitor. The disambiguator.} \\
\textbf{\code{Quantity}} & \textbf{Units on the shelf. Copied out.} \\
\code{Latest Cost} & Most recent purchase cost. Unused by the tools. \\
\textbf{\code{Average Cost}} & \textbf{Averaged purchase cost. Copied out.} \\
\code{Sale Price} & Retail price. Unused by the tools. \\
\bottomrule
\end{tabular}
\end{center}

Note that the tools copy \code{Average Cost}, not \code{Latest Cost} and not
\code{Sale Price}. The README describes this tool as updating ``quantities,
costs''; that is accurate. It does not touch the customer-facing price at all.

\subsection{The Shopify export}

Fifteen columns in the sample, following the Shopify product-export layout. The
ones that matter:

\begin{center}
\begin{tabular}{@{}ll@{}}
\toprule
\textbf{Column} & \textbf{Role} \\
\midrule
\code{Handle} & Groups the rows of one product together. \\
\textbf{\code{Title}} & \textbf{Product name. Join key for single-variant products.} \\
\textbf{\code{Body (HTML)}} & \textbf{The description, as raw HTML. Holds the dimensions table.} \\
\code{Option1 Name} & \code{Title} for single-variant products, else e.g. \code{Size}. \\
\textbf{\code{Option1 Value}} & \textbf{\code{Default Title}, or the variant name. Join key for variants.} \\
\code{Variant Inventory Qty} & What the storefront currently believes. Read by nothing. \\
\code{Cost per item} & What the storefront currently believes. Read by nothing. \\
\bottomrule
\end{tabular}
\end{center}

\begin{keyidea}{The \code{Default Title} sentinel}
Shopify gives every product at least one variant. A product with no real variants
gets one anyway, and Shopify labels it \code{Option1 Value = Default Title}. So
that literal string is how a program tells ``this product has no variants, its
identity is in \code{Title}'' apart from ``this is variant N of a product, its
identity is in \code{Option1 Value}''. Every tool in this repository branches on
that string.
\end{keyidea}

Observe the consequence in the sample: for the second variant row of product
\code{22210}, the \code{Title} cell is empty. Shopify writes the product-level
fields only on the first row of each \code{Handle} group. That is exactly why
variants cannot be matched on \code{Title} and need \code{Option1 Value}.

\subsection{What the sample data actually is}

\begin{honest}{The ``11,000-line Shopify export'' is not in this repository}
The README states, correctly and in the past tense, that the original repo held
the client's real data: roughly 1,600 stock rows and an 11,000-line Shopify
export, and that it was removed. Confirmed against what is actually here:

\begin{itemize}
  \item \file{sample\_shopify\_stock.csv} is \textbf{117 physical lines, which
  pandas reads as 6 product rows} (the HTML descriptions span many lines each).
  \item \file{sample\_stock\_list.csv} is \textbf{8 stock rows}.
\end{itemize}

So the figures ``1,600'' and ``11,000'' describe data that no longer exists in
the repository, and \textbf{nothing in the repo can corroborate them}. They rest
entirely on the owner's recollection of the 2023 engagement. Git history was
checked: the real data was \textbf{never committed}, not even in the initial
commit, so it is not recoverable from history either. That is good for
confidentiality and it means the numbers are unverifiable by anyone, including
the owner, without the original client files.

This is not an accusation that the numbers are wrong. It is a warning that they
are the one class of claim on this project that cannot be defended by pointing at
the code. If an interviewer asks ``how do you know it was 11,000 lines'', the
honest answer is ``I worked on the file in 2023; it is not in the repo because it
was client-internal''. Do not improve the number, and do not claim the sample
demonstrates scale. It demonstrates correctness, at six rows.
\end{honest}

The sample is well built for its actual job. It is deliberately adversarial in
two places: \code{6205} and \code{6304 2RS} each have a competing \code{OTHER}
brand row with different figures, so a tool that forgets to filter by brand
produces visibly wrong output. That is a test fixture in all but name.

\section{Tool 1: the quantity and pricing updater}
\label{sec:matching}

\file{quantity\_pricing\_updater/update\_quantity\_pricing.py}. 69 lines, no
functions, top-to-bottom script.

\subsection{Entry point}

\begin{lstlisting}[language=Python,caption={Argument handling}]
if len(sys.argv) == 3:
    stock_list_path, shopify_stock_path = sys.argv[1], sys.argv[2]
else:
    from tkinter.filedialog import askopenfilename
    ...
    while stock_list_path == '':
        stock_list_path = askopenfilename(filetypes=[('Stock List CSV', '*.csv')])
\end{lstlisting}

Two paths on the command line, or a file picker for each. The \code{while} loop
means cancelling the dialog re-opens it forever: there is no way to abort except
killing the process.

\begin{jargon}{pandas and the DataFrame}
\code{pandas} is the standard Python library for tabular data. Its central object
is the \textbf{DataFrame}: a table in memory with named columns, loaded from a CSV
in one call. \code{df[df['Brand'] == 'HI-TEC']} produces a new DataFrame of only
the rows where that is true. This filtering style is what the whole project is
built on.
\end{jargon}

\subsection{The matching problem}

\begin{keyidea}{Two tables with no shared key}
The stock list has no Shopify product ID. The Shopify export has no stock-system
item number. There is no column that means the same thing in both files and can
be trusted to be unique. All that exists is a bearing designation that appears in
both, as text: \code{6205} in the stock list's \code{Item}, and \code{6205} in
the Shopify export's \code{Title}. So the join is on text, and the text is not
unique, because \code{6205} appears once per brand.
\end{keyidea}

The disambiguator is \code{Brand}. Only the \code{HI-TEC} row is the one to
trust. Filter first, then join:

\begin{figure}[H]
\centering
\resizebox{\textwidth}{!}{
\begin{tikzpicture}[node distance=7mm and 14mm]
  \node[store] (sl) {stock list\\8 rows};
  \node[box, right=of sl] (filt) {filter\\\code{Brand == 'HI-TEC'}};
  \node[store, right=of filt] (hitec) {6 trusted rows\\\code{Item, Quantity,}\\\code{Average Cost}};

  \node[store, below=16mm of sl] (ss) {Shopify export\\6 rows};
  \node[dec, right=of ss] (split) {\code{Option1 Value}\\\code{== 'Default}\\\code{Title'}?};
  \node[box, above right=2mm and 12mm of split] (defl) {default-title view\\join on \code{Title}};
  \node[box, below right=2mm and 12mm of split] (nond) {variant view\\join on \code{Option1 Value}};
  \node[box, right=22mm of split, xshift=28mm] (cat) {\code{concat}};
  \node[store, right=of cat] (out) {\code{updated\_*.csv}\\+ \code{Updated Quantity}\\+ \code{Updated Cost}};

  \draw[flow] (sl) -- (filt);
  \draw[flow] (filt) -- (hitec);
  \draw[flow] (ss) -- (split);
  \draw[flow] (split) -- node[lbl,above,sloped]{yes} (defl);
  \draw[flow] (split) -- node[lbl,below,sloped]{no} (nond);
  \draw[flow] (defl) -| (cat);
  \draw[flow] (nond) -| (cat);
  \draw[flow] (cat) -- (out);
  \draw[dflow] (hitec) -- node[lbl,right]{lookup\\source} (defl);
\end{tikzpicture}}
\caption{The data flow. The brand filter happens once, up front. The split on the
\code{Default Title} sentinel decides which column is the join key.}
\end{figure}

\subsection{The matching loop, line by line}

\begin{lstlisting}[language=Python,caption={The default-title branch, verbatim}]
default_title_shopify_stock["Updated Quantity"] = 0
default_title_shopify_stock["Updated Cost"] = 0

title_shopify_stock = default_title_shopify_stock["Title"].tolist()
for x in title_shopify_stock:
    if x in stock_list[stock_list['Brand'] == 'HI-TEC']["Item"].tolist():
        print('original', x)
        ...
        default_title_shopify_stock.loc[
            default_title_shopify_stock["Title"] == x,
            ['Updated Quantity', 'Updated Cost']] = stock_list.loc[
            (stock_list["Item"] == x) & (stock_list['Brand'] == 'HI-TEC'),
            ['Quantity', 'Average Cost']].values[0]
\end{lstlisting}

Read carefully, this is doing the following per Shopify row:

\begin{pseudocode}{What one iteration costs}
\Step{1. Re-filter the ENTIRE stock list to \code{Brand == 'HI-TEC'}.}
\Step{2. Convert that filtered column to a fresh Python list.}
\Step{3. Linear-scan the list for the title.}
\Step{4. If found, filter the entire stock list AGAIN (twice more, for the two}
\Indent{\code{.loc} lookups), take \code{.values[0]}.}
\Step{5. Scan the whole Shopify view again to find the rows to write to.}
\end{pseudocode}

\begin{honest}{The matching is quadratic, and needlessly}
Steps 1 to 3 rebuild the same filtered list on every single iteration. Nothing in
the loop changes the stock list, so this is the same work re-done $n$ times. With
$n$ Shopify rows and $m$ stock rows the loop is $O(n \times m)$ with a large
constant factor, when the entire operation is one \code{pandas.merge} on
\code{Item} against a pre-filtered frame: $O(n + m)$, and about four lines. The
README says this itself and is right to. On six rows it is invisible. On the
11,000-row export it is the difference between instant and a coffee break.

Worth being precise about the defence: this is not a subtle bug, it is a
straightforwardly slow implementation of a correct idea. The idea (filter by
brand, then join) is the part that took the thinking, and it is right.
\end{honest}

\begin{honest}{Writing into a filtered view: correct here, by luck of version}
\code{default\_title\_shopify\_stock} is the result of \code{shopify\_stock[mask]}.
In pandas that is a \textbf{view-or-copy of ambiguous identity}, and assigning
into it raises \code{SettingWithCopyWarning}. Running the tool prints exactly
that warning.

It works today because the script never relies on the write propagating back to
the parent frame: it writes into the two views, then \code{concat}s the two views
together and saves that. So the ambiguity is dodged rather than resolved.

\textbf{This is why \file{requirements.txt} pins \code{pandas>=2.0,<3.0}.} Under
pandas 3.0's copy-on-write semantics the assignment target becomes unambiguously
a copy, and code of this shape loses writes silently rather than warning. The pin
is load-bearing and the README knows it. Removing the upper bound is not a
maintenance chore; it is a correctness change that needs the loops rewritten
first.
\end{honest}

\subsection{Verified behaviour}

The tool was run against the shipped sample.

\begin{center}
\begin{tabular}{@{}llrrl@{}}
\toprule
\textbf{Title} & \textbf{Option1 Value} & \textbf{Updated Qty} & \textbf{Updated Cost} & \textbf{Note} \\
\midrule
\code{6205} & \code{Default Title} & 55 & 50 & HI-TEC row, not \code{OTHER}'s 3 / 45 \\
\code{6206} & \code{Default Title} & 30 & 64 & \\
\code{1205} & \code{Default Title} & 28 & 85 & \\
\code{6304 2RS} & \code{Default Title} & 70 & 41 & HI-TEC row, not \code{OTHER}'s 5 / 38 \\
\code{22210} & \code{22210 K} & 15 & 172 & matched via \code{Option1 Value} \\
(empty) & \code{22210 E} & 10 & 185 & matched despite empty \code{Title} \\
\bottomrule
\end{tabular}
\end{center}

\begin{keyidea}{Confirmed: the brand filter earns its place}
Both trap rows resolve to the HI-TEC figures. Both variant rows resolve through
\code{Option1 Value}, including the one whose \code{Title} is blank. The README's
central claim about this tool is true, and I checked it by running it rather than
reading it.
\end{keyidea}

The output also confirms a detail the README does not mention: rows that match
nothing keep \code{Updated Quantity = 0} rather than blank. A commented-out block
in the source shows the author considered converting those zeros to \code{NaN}
and did not:

\begin{lstlisting}[language=Python,caption={Dead code, left in the GUI version}]
# # Replace values with 0 to NaN in Updated Quantity and Updated Cost
# shopify_stock['Updated Quantity'] = shopify_stock['Updated Quantity'].replace(0, pandas.np.nan)
\end{lstlisting}

\begin{honest}{A zero that means ``no data'' is a hazard in this workflow}
An unmatched product gets \code{Updated Quantity = 0}, which is indistinguishable
from a product genuinely out of stock. Since the output spreadsheet is re-imported
into a live shop, a human who trusts that column without checking would zero the
inventory of every product that failed to match. The review step is what catches
this, which is precisely why the review step exists. The commented-out
\code{NaN} line would have made the distinction visible. It is also broken as
written: \code{pandas.np} was removed from pandas in version 2.0, so that line
would raise \code{AttributeError} if uncommented under the pinned version.
\end{honest}

\subsection{Output path construction}

\begin{lstlisting}[language=Python]
updated_file_path = shopify_stock_path.split('/')
updated_file_path[-1] = 'updated_' + updated_file_path[-1]
updated_file_path = '/'.join(updated_file_path)
\end{lstlisting}

String surgery on the path, splitting on a hard-coded forward slash. On Windows,
where the client's packaged executable ran and where paths use backslashes, this
is saved only by \code{askopenfilename} returning forward-slash paths regardless
of platform. It is fragile by accident rather than by design; \code{os.path} or
\code{pathlib} does this correctly in one call. The CLI version also silently
overwrites an existing \code{updated\_} file. Only the GUI guards against that.

\section{Tool 2: the dimensions extractor}

\file{dimensions\_extractor/dimensions\_extractor.py}. 83 lines.

\subsection{The problem: data locked inside prose}

Each product's \code{Body (HTML)} contains a dimensions table as raw markup:

\begin{lstlisting}[language=HTML,caption={A real \code{Body (HTML)} cell from the sample}]
<p>Product Category: Ball Bearings</p>
<table width="100%">
<tbody>
<tr><td colspan="3">Dimensions(mm)</td></tr>
<tr><td>d</td><td>Bore Diameter</td><td>25</td></tr>
<tr><td>D</td><td>Outside Diameter</td><td>52</td></tr>
<tr><td>T</td><td>Width</td><td>15</td></tr>
</tbody>
</table>
\end{lstlisting}

(Shown here folded onto single lines for width. In the file each \code{<td>} is
on its own line, and that fact is load-bearing, as the next section explains.)

\begin{jargon}{HTML tag}
Text inside angle brackets that marks up a document's structure.
\code{<td>25</td>} means ``a table cell containing 25''. The \code{<td>} and
\code{</td>} are the tags; \code{25} is the content. A browser reads the tags and
draws a table. A spreadsheet column just holds the whole thing as one long string.
\end{jargon}

The customer sees a neat table. The storefront cannot filter on ``bore diameter
between 20 and 30'' because as far as Shopify is concerned this is one opaque
blob of description text. Getting those three numbers into their own columns is
what makes them usable as metafields.

\begin{jargon}{Metafield}
A custom field attached to a Shopify product beyond the built-in ones. Once bore
diameter is a metafield rather than a sentence in the description, the theme can
offer a filter on it. This is the business reason the tool exists.
\end{jargon}

\subsection{The strategy: flatten and read line pairs}

\begin{figure}[H]
\centering
\resizebox{\textwidth}{!}{
\begin{tikzpicture}[node distance=6mm and 10mm]
  \node[store] (raw) {\code{Body (HTML)}\\raw markup string};
  \node[box, right=of raw] (low) {\code{.lower()}};
  \node[box, right=of low] (strip) {\code{re.sub('<.*?>', '')}\\strip every tag};
  \node[box, right=of strip] (split) {split on newline,\\drop empty lines};
  \node[store, right=of split] (lines) {flat list\\of text lines};
  \node[box, below=12mm of lines] (scan) {find index of\\label line};
  \node[box, left=of scan] (next) {take the\\NEXT line};
  \node[dec, left=of next] (num) {numeric?};
  \node[store, left=of num] (col) {write into\\column};
  \node[extbox, below=8mm of num] (skip) {skip: leave\\column empty};

  \draw[flow] (raw) -- (low);
  \draw[flow] (low) -- (strip);
  \draw[flow] (strip) -- (split);
  \draw[flow] (split) -- (lines);
  \draw[flow] (lines) -- (scan);
  \draw[flow] (scan) -- (next);
  \draw[flow] (next) -- (num);
  \draw[flow] (num) -- node[lbl,above]{yes} (col);
  \draw[flow] (num) -- node[lbl,left]{no} (skip);
\end{tikzpicture}}
\caption{The extraction pipeline. Note what survives: the table becomes a flat
list of lines, and all structure is gone except adjacency.}
\end{figure}

\begin{keyidea}{The core assumption}
After the tags are stripped, \code{<td>Bore Diameter</td>} and \code{<td>25</td>}
were on separate lines, so they become two consecutive entries in the list:
\code{['bore diameter', '25']}. So: \textbf{find the label, take the very next
line, and that is the value}. The whole extractor rests on this one assumption.
\end{keyidea}

\begin{lstlisting}[language=Python,caption={The tag stripper}]
def remove_html_tags(text_to_be_cleaned):
    """Remove html tags from a string"""
    clean = re.compile('<.*?>')
    return re.sub(clean, '', text_to_be_cleaned)
\end{lstlisting}

\begin{jargon}{Regular expression (regex)}
A pattern language for text. \code{<.*?>} means: a \code{<}, then any characters,
then a \code{>}. The \code{?} makes it \textbf{non-greedy}, so it stops at the
first \code{>} rather than running to the last one on the line. Without the
\code{?}, the pattern would eat \code{<td>25</td>} whole and the number would be
destroyed along with the tags. That single character is why this function works.
\end{jargon}

\subsection{The scanning loop and its three branches}

\begin{lstlisting}[language=Python,caption={The heart of the extractor}]
values_to_check = ['outside diameter', 'bore diameter', 'width']
for x in df.iterrows():
    temp = str(x[1]['Body (HTML)']).lower()
    text = [x for x in remove_html_tags(temp).split('\n') if x != '']
    for value in values_to_check:
        if value in text:                          # branch 1: exact line match
            index = text.index(value) + 1
            if check_int_float(text[index]):
                df[value][x[0]] = text[index]
        elif value == 'width':                     # branch 2: the 'total width' fallback
            if 'total width' in text:
                index = text.index('total width') + 1
                ...
        else:                                      # branch 3: substring scan
            for temp in text:
                if value in temp:
                    index = text.index(temp) + 1
                    ...
\end{lstlisting}

Three ways to find a label, tried in order:

\begin{enumerate}
  \item \textbf{Exact match.} The line is exactly \code{bore diameter}. This is
  the normal path and the only one the sample exercises.
  \item \textbf{The \code{total width} fallback.} Some products label the field
  \code{total width} rather than \code{width}. This branch is real
  domain knowledge, learned from the client's actual catalogue, and it is the
  most interesting line in the file: it is the fingerprint of someone who looked
  at the data rather than assuming it.
  \item \textbf{Substring scan.} If a line merely \emph{contains} the label, use
  it. This catches \code{bore diameter (mm)} and similar.
\end{enumerate}

\begin{honest}{Branch 2 shadows branch 3 for \code{width}, permanently}
The \code{elif value == 'width'} sits between the exact match and the substring
scan. So for \code{width}, and only for \code{width}, branch 3 is unreachable: if
the exact line is absent, control goes to branch 2 and stops there whether or not
it finds \code{total width}. A line reading \code{overall width (mm)} would be
found for \code{bore diameter}-style labels but never for \code{width}. The
ordering is almost certainly unintentional.
\end{honest}

\begin{honest}{\code{text.index(temp)} is the wrong lookup in branch 3}
\code{for temp in text} already iterates the lines, but the code then calls
\code{text.index(temp)} to ask where that line is. \code{list.index} returns the
\textbf{first} occurrence. If the same line text appears twice in a description,
branch 3 reads the value following the first one regardless of which one it
matched. \code{enumerate} gives the correct index for free. The same pattern
recurs in the variant detector.
\end{honest}

\begin{honest}{The loop variable \code{temp} is reused for two different things}
\code{temp} holds the lowercased HTML string, and then branch 3 rebinds it as the
per-line loop variable, destroying the original inside that iteration. It happens
to be harmless because the HTML string is not read again in that pass, but it is
the kind of accident that stops being harmless the moment a line is added.
The outer loop variable \code{x} is likewise shadowed by the list comprehension
that builds \code{text}.
\end{honest}

\subsection{Verified behaviour}

Run against the sample, the extractor recovered every dimension:

\begin{center}
\begin{tabular}{@{}llrrr@{}}
\toprule
\textbf{Title} & \textbf{Option1 Value} & \textbf{bore} & \textbf{outside} & \textbf{width} \\
\midrule
\code{6205} & \code{Default Title} & 25 & 52 & 15 \\
\code{6206} & \code{Default Title} & 30 & 62 & 16 \\
\code{22210} & \code{22210 K} & 50 & 90 & 23 \\
\code{1205} & \code{Default Title} & 25 & 52 & 15 \\
\code{6304 2RS} & \code{Default Title} & 20 & 52 & 15 \\
\bottomrule
\end{tabular}
\end{center}

The README's claim that it ``recovers all three dimensions from the sample HTML''
is confirmed. But count the rows.

\begin{honest}{The output has five rows. The input had six.}
\begin{lstlisting}[language=Python]
df = df.dropna(subset=['Body (HTML)'])
\end{lstlisting}
The \code{22210 E} variant has no description of its own (Shopify writes
product-level fields only on the first row of a \code{Handle} group), so its
\code{Body (HTML)} is empty, so \code{dropna} \textbf{deletes the row from the
output file entirely}.

This is a real defect and the README does not mention it. The tool's contract, as
described, is ``adds three columns''. Its actual behaviour is ``adds three columns
and silently drops every variant row that lacks its own description'', which on a
catalogue of multi-variant products is potentially most of the rows. The intent
was obviously to skip such rows when scanning; the effect is to remove them from
the deliverable.

How bad is it in practice? The output is re-imported to Shopify, and a Shopify
product import updates the rows present rather than deleting absent ones, so this
most likely does not destroy live variants. It does mean those variants never
receive their dimension columns, quietly. Marking this \textbf{inferred}: I have
not verified Shopify's import behaviour against a live store, and doing so was out
of scope. Either way, \code{df.loc[df['Body (HTML)'].notna()]} as the loop's
iteration source, rather than \code{dropna} on the frame itself, is the one-line
fix that keeps every row in the output.
\end{honest}

\begin{honest}{The two versions disagree on the column's type}
The CLI writes \code{df[value][x[0]] = text[index]}, storing the \textbf{string}
\code{'25'}. The GUI writes \code{float(text[index])}, storing the \textbf{number}
\code{25.0}. Same tool, same column, two types, depending on which binary the
client double-clicked. For a filterable metafield that difference decides whether
sorting is numeric or lexicographic, which is the difference between
$9 < 10$ and \code{'10' < '9'}.
\end{honest}

Both versions use \textbf{chained assignment}, \code{df[value][x[0]] = ...},
which selects a column and then indexes into it. This is the pattern the pandas
documentation warns about most loudly, and the second reason for the
\code{pandas<3} pin.

\section{Tool 3: the GUI, and a third operation}

\file{stock\_updater\_gui/shopify\_stock\_updater.py}. 354 lines. This is what
the client actually received.

\begin{jargon}{tkinter}
Python's built-in GUI toolkit. It ships with Python, needs no installation on
Windows, and produces a native desktop window. That last property is the entire
reason it was chosen: the deliverable had to be a thing a non-technical person
double-clicks.
\end{jargon}

\begin{figure}[H]
\centering
\begin{tikzpicture}[node distance=5mm and 10mm]
  \node[box, minimum width=52mm] (win) {\code{root} window, 900x800, non-resizable};
  \node[box, below=of win, minimum width=52mm] (b1) {Button: Update Quantity and Pricing};
  \node[box, below=of b1, minimum width=52mm] (b2) {Button: Extract Data from HTML};
  \node[box, below=of b2, minimum width=52mm] (b3) {Button: Variant Detection};
  \node[box, below=of b3, minimum width=52mm, fill=white] (err) {\code{error}: read-only Text widget\\(status and error display)};

  \node[extbox, right=22mm of b1] (f1) {\code{update\_quantity\_and\_pricing()}};
  \node[extbox, right=22mm of b2] (f2) {\code{extract\_data\_from\_html()}};
  \node[extbox, right=22mm of b3] (f3) {\code{variants\_detection()}};
  \node[box, right=22mm of err] (sv) {\code{save\_data()}};

  \draw[flow] (b1) -- (f1); \draw[flow] (b2) -- (f2); \draw[flow] (b3) -- (f3);
  \draw[flow] (f1) -- (sv); \draw[flow] (f2) -- (sv); \draw[flow] (f3) -- (sv);
  \draw[dflow] (sv) -- node[lbl,above]{writes status} (err);
\end{tikzpicture}
\caption{The GUI's structure. Three buttons, three handlers, one shared save
routine, one status widget. Verified against the shipped screenshot
\file{docs/screenshots/app.png}, which shows exactly this layout.}
\end{figure}

\subsection{Additions over the CLI tools}

\begin{itemize}
  \item \textbf{XLSX input and output.} \code{open\_file} inspects the chosen
  extension and returns a boolean alongside the path, so callers pick
  \code{read\_excel} or \code{read\_csv}. Output is always XLSX.
  \item \textbf{Distinct output prefixes.} \code{updated\_p\_} for pricing,
  \code{updated\_h\_} for HTML dimensions, \code{updated\_v\_} for variants, so the
  three operations do not clobber each other.
  \item \textbf{Non-destructive saves.} See \code{save\_data} below.
  \item \textbf{Hover feedback.} \code{on\_enter} and \code{on\_leave} swap the
  button background. Small, but it is the kind of thing that makes software feel
  alive to a non-technical user.
\end{itemize}

\begin{honest}{\code{open\_file}'s retry loop does not loop}
\begin{lstlisting}[language=Python]
filepath = ''
while filepath == '':
    ...
    filepath = askopenfilename(...)
    if filepath.split('.')[-1] == 'xlsx':
        return filepath, True
    else:
        return filepath, False
\end{lstlisting}
Both branches of the \code{if} return, so the \code{while} body can never run
twice. If the user cancels the dialog, \code{askopenfilename} returns the empty
string, the \code{else} branch returns \code{('', False)}, and the caller hands
that to \code{read\_csv}, which raises and takes the handler down with an
unhandled exception. The window survives (tkinter catches it and prints to the
console) but the user sees the prompt text and no explanation, because
nothing writes to the \code{error} widget on failure. The loop was clearly meant
to re-prompt; the returns defeat it.

There is a second, quieter bug in the same three lines: the extension test is
\code{filepath.split('.')[-1] == 'xlsx'}. Cancel returns \code{''}, whose split
is \code{['']}, so the \code{else} fires. But any file whose name has no dot, or
a path containing a dot in a directory name, is classified by the wrong segment.
\end{honest}

\subsection{\code{save\_data} and the non-overwrite counter}

\begin{figure}[H]
\centering
\begin{tikzpicture}[node distance=8mm and 14mm]
  \node[box] (start) {\code{save\_data(path)}};
  \node[dec, below=of start] (try) {\code{read\_excel(path)}\\succeeds?};
  \node[box, right=18mm of try] (bump) {\code{counter += 1}\\path becomes\\\code{(N)name.xlsx}};
  \node[box, below=14mm of try] (write) {\code{to\_excel(path)}\\status: green message};

  \draw[flow] (start) -- (try);
  \draw[flow] (try) -- node[lbl,above]{yes: file exists} (bump);
  \draw[flow] (bump) -- ++(0,16mm) -| (try);
  \draw[flow] (try) -- node[lbl,left]{\code{FileNotFoundError}} (write);
\end{tikzpicture}
\caption{The non-overwrite loop. ``Does this file exist'' is answered by trying
to parse it as a whole spreadsheet.}
\end{figure}

\begin{keyidea}{The intent is right}
The client re-runs these tools repeatedly on the same folder. Silently overwriting
yesterday's reviewed output would be a genuinely destructive default. Bumping to
\code{(1)name.xlsx}, \code{(2)name.xlsx} and so on mirrors what a browser does for
duplicate downloads, which is exactly the right idiom for this audience.
\end{keyidea}

\begin{honest}{The existence check is a full spreadsheet parse}
\begin{lstlisting}[language=Python]
try:
    pandas.read_excel(updated_file_path)   # existence check
    counter += 1
    ...
except FileNotFoundError:
    df.to_excel(updated_file_path, index=False)
\end{lstlisting}
To learn whether a path exists, the code \textbf{reads and parses the entire
workbook} and throws the result away. On the real 11,000-row export, saving the
tenth run means parsing nine full spreadsheets first. \code{os.path.exists(path)}
is the one-call answer.

Worse, it catches only \code{FileNotFoundError}. A file that exists but is not a
readable workbook (a partially written file, a permissions error, or the very
common case of \textbf{the file being open in Excel}) raises something else:
\code{PermissionError}, \code{BadZipFile}, \code{ValueError}. None are caught, so
the handler dies and no status ever reaches the user. Given that the audience is a
client who lives in Excel and will absolutely have the previous output open while
running the tool again, this is the failure mode most likely to have actually bitten
in practice.
\end{honest}

\subsection{Variant detection}

The third operation, present only in the GUI. It parses bearing designation
suffixes into human-readable columns.

\begin{keyidea}{The domain knowledge}
A bearing designation like \code{22210 EK} is not a name, it is an encoding. The
digits give the series and size; the trailing letters describe construction. This
function is a lookup table for that code, and it is the only place in the
repository where the author's understanding of the client's product domain is
written down.
\end{keyidea}

\begin{center}
\begin{tabular}{@{}llll@{}}
\toprule
\textbf{Column} & \textbf{Token} & \textbf{Mapped to} & \\
\midrule
\code{Bore} & \code{k} & \code{Tapered} & \\
\midrule
\code{Cage} & \code{e} & \code{Steel} & \\
 & \code{f} & \code{Fibre} & \\
 & \code{m} & \code{Brass} & \\
 & \code{ecp} & \code{Polyamide} & \\
 & \code{ca} & \code{Brass} & \\
 & \code{mb} & \code{Steel} & \\
 & \code{cc} & \code{Brass} & \\
\midrule
\code{Seal Type} & \code{2rs} & \code{2RS} & \\
 & \code{rs} & \code{rs} & (unchanged) \\
 & \code{z} & \code{z} & (unchanged) \\
 & \code{2z} & \code{2z} & (unchanged) \\
\bottomrule
\end{tabular}
\end{center}

Plus two rules that run after the token scan:

\begin{pseudocode}{The series-2xxxx override}
\Step{if the variant's last token is numeric:}
\Indent{Bore := 'Plain', Cage := 'Steel'}
\Indent{if first token starts with '2' and has length 5: Cage := 'Brass'}
\Step{else, after scanning tokens:}
\Indent{if 'e' or 'ca' or 'mb' in tokens:}
\Indent{\ \ if numeric token starts with '2' and has length 5: Cage := 'Brass'}
\Indent{if 'cc' in tokens:}
\Indent{\ \ if numeric token starts with '2' and has length 5: Cage := 'Steel'}
\end{pseudocode}

Traced against the sample:

\begin{center}
\begin{tabular}{@{}lllll@{}}
\toprule
\textbf{Variant} & \textbf{Tokens} & \textbf{Bore} & \textbf{Cage} & \textbf{Seal} \\
\midrule
\code{22210 K} & \code{['22210','k']} & \code{Tapered} & (empty) & (empty) \\
\code{22210 E} & \code{['22210','e']} & (empty) & \code{Brass} & (empty) \\
\bottomrule
\end{tabular}
\end{center}

For \code{22210 E}: the token scan sets \code{Cage = 'Steel'} from the \code{e}
lookup, then the override fires (first token \code{22210} starts with \code{2},
length 5) and rewrites it to \code{Brass}. The table lookup is deliberately
overruled by the series rule. That is real domain knowledge encoded as a
post-pass, and it is the most defensible thing in this function.

\begin{honest}{\code{sealtype\_possible\_variants\_name} is three-quarters a no-op}
\begin{lstlisting}[language=Python]
sealtype_possible_variants = ['2rs', 'rs', 'z', '2z']
sealtype_possible_variants_name = ['2RS', 'rs', 'z', '2z']
\end{lstlisting}
Only \code{2rs} is translated, and only into its own uppercase form. The other
three map to themselves. The README says variant detection parses suffixes ``into
human-readable \code{Bore}, \code{Cage}, and \code{Seal Type} columns''. For
\code{Bore} and \code{Cage} that is true and useful. For \code{Seal Type} it is
not: the column receives the raw code back, unchanged and mostly lowercase. Either
the mapping was never finished, or the seal codes were considered self-explanatory
and the parallel-list structure was kept for symmetry. The code does not say which.
\end{honest}

\begin{honest}{Parallel lists are the wrong data structure here}
Every lookup is \code{name\_list[code\_list.index(token)]}: two lists kept in
sync by hand, joined by a positional index. A single \code{dict} makes the
correspondence impossible to break, halves the code, and makes
\code{'m': 'Brass'} readable as a fact rather than as ``the third element of one
list goes with the third element of another''. As written, inserting a code into
one list without the other silently shifts every mapping after it.
\end{honest}

\begin{honest}{\code{x} is set from the wrong index, and defaults to 0}
\begin{lstlisting}[language=Python]
x = 0
for value in item_individual:
    if check_int_float(value):
        x = item_individual.index(value)
\end{lstlisting}
\code{x} is meant to be the position of the numeric token, and it is later used as
\code{item\_individual[x]} in the series override. Two problems: \code{index()}
finds the first occurrence rather than the current one, and if there is no numeric
token at all, \code{x} silently stays \code{0} and the override tests token 0
anyway. It works on the sample because the numeric token \emph{is} token 0. On
\code{K 22210} it would still work by accident; on a designation with two numeric
tokens it would read the wrong one.
\end{honest}

\begin{honest}{The variant detector ignores \code{Brand} entirely}
Tool 1's central insight is that a designation is ambiguous without the brand.
The variant detector then parses designations with no brand filter at all. This
is defensible, since a suffix code means the same thing whoever made the bearing,
but it is worth noticing that the repository's own hard-won rule is not applied
here, and no comment says why.
\end{honest}

\section{The GUI/CLI divergence}

\begin{keyidea}{The GUI is a copy, not a wrapper}
The README says the GUI ``bundles both tools above''. Read the source and it does
not import them; it contains its own near-identical copies of both. The README's
``What I'd do differently'' names this: the logic should live in one shared module
with two frontends. It is right, and this section is what that duplication cost.
\end{keyidea}

Duplication means the two copies drift, and they have. The type disagreement in
the dimensions column is one instance. Here is the serious one.

\subsection{A verified bug: the GUI's quantity matcher matches nothing}

The GUI's \code{update\_quantity\_and\_pricing} has one line its CLI sibling does
not:

\begin{lstlisting}[language=Python,caption={GUI only. This line breaks the match.}]
def convert_int_float(string):
    """Check if a string is an integer or a float"""
    try:
        num = float(string)
        return num
    except ValueError:
        pass          # <- returns None implicitly

stock_list['Item'] = stock_list['Item'].apply(convert_int_float)
\end{lstlisting}

The intent is evidently type normalisation: force the \code{Item} column to a
consistent numeric type before matching. What it actually does, traced against the
shipped sample:

\begin{center}
\begin{tabular}{@{}lll@{}}
\toprule
\textbf{Stage} & \textbf{\code{Item} column} & \\
\midrule
As read from CSV & \code{['6205','6205','6206','22210 K','22210 E',...]} & all strings \\
After \code{convert\_int\_float} & \code{[6205.0, 6205.0, 6206.0, nan, nan, ...]} & floats and \code{nan} \\
\bottomrule
\end{tabular}
\end{center}

Two things happened:

\begin{enumerate}
  \item \code{'6205'} became the float \code{6205.0}. The Shopify \code{Title}
  it must match is still the \textbf{string} \code{'6205'}, and
  \code{'6205' in [6205.0]} is \code{False}. Python does not compare a string
  equal to a float.
  \item \code{'22210 K'} is not parseable as a float, so \code{convert\_int\_float}
  falls through to \code{pass} and \textbf{returns \code{None}}, which pandas
  stores as \code{nan}. Every alphanumeric designation, which is to say every
  variant, is destroyed outright.
\end{enumerate}

I ran the GUI's exact matching logic against the shipped sample. The result:

\begin{center}
\begin{tabular}{@{}ll@{}}
\toprule
\textbf{Branch} & \textbf{Matches found} \\
\midrule
Default-title rows (join on \code{Title}) & \textbf{0} of 4 \\
Variant rows (join on \code{Option1 Value}) & \textbf{0} of 2 \\
\bottomrule
\end{tabular}
\end{center}

\begin{honest}{The GUI's flagship button produces an all-zero column on this data}
The CLI version matches all six rows correctly. The GUI version, given the same
two files, matches none, writes \code{Updated Quantity = 0} and
\code{Updated Cost = 0} for every product, reports ``Quantity and Pricing has been
updated!'' in green, and saves the file. It fails silently and then congratulates
itself.

And this is the version the client received.

\textbf{Scope of the claim, stated precisely.} What is verified is that the GUI's
matching logic finds zero matches \emph{on the sample data shipped in this
repository}. Whether it failed on the client's real 2023 files is
\textbf{inferred and unresolvable from this repo}, because it depends on the dtypes
pandas inferred from data that no longer exists. The mechanism is dtype-dependent
in a specific way:

\begin{itemize}
  \item If the real stock list's \code{Item} column was \textbf{purely numeric},
  pandas would read it as \code{int64} or \code{float64}, the real Shopify
  \code{Title} column likewise, and \code{6205.0 == 6205.0} would match. The line
  would then be harmless, or even the thing that made the types agree.
  \item If either column contained \textbf{any} alphanumeric designation, and the
  sample says the real data did (\code{22210 K}, \code{6304 2RS} are exactly the
  kind of thing this catalogue is full of), that column is \code{object} dtype,
  its entries are strings, and the match collapses as it does here.
\end{itemize}

The second is far more likely, which raises an uncomfortable and genuinely open
question: either the GUI's pricing button was broken in the client's hands, or the
real data had a shape the sample does not reproduce. The repository cannot settle
it. \textbf{This is the single most important thing to know about this project
before discussing it}, and it is not in the README.

The honest position, and a strong one: the CLI implementation is correct and
verifiably so; the GUI copy carries a well-understood dtype bug; the bug is the
direct, predictable consequence of the duplication the README already identifies
as the project's main design regret. That is a real lesson with a real receipt
attached, and it is worth more than a clean bill of health.

The fix is two lines: drop the \code{convert\_int\_float} normalisation and
compare as strings on both sides, exactly as the working CLI version does.
\end{honest}

\begin{honest}{\code{convert\_int\_float} has no explicit return, and a copied docstring}
\code{check\_int\_float} returns a bool. \code{convert\_int\_float} was written by
copying it, and the docstring came along unedited: both say ``Check if a string is
an integer or a float'', though the second converts rather than checks. Its
failure path is a bare \code{pass}, so it returns \code{None} implicitly. An
explicit \code{return None}, or better a \code{return string} that leaves
unparseable values alone, would have made the \code{nan} destruction visible at
the point it was written.
\end{honest}

\section{Packaging and dependencies}

\subsection{PyInstaller}

\begin{jargon}{PyInstaller}
A tool that bundles a Python script, the Python interpreter itself, and every
imported library into a single executable. The recipient runs it without
installing Python. A \file{.spec} file is the build recipe.
\end{jargon}

Both CLI tools have a \file{.spec}. They are byte-identical apart from the script
name and the output name, and both are stock PyInstaller output:
\code{console=True}, \code{upx=True}, one-file, no hidden imports, no bundled data.

\begin{honest}{The GUI, the one thing the client actually ran, has no \file{.spec}}
The README states this plainly in its ``Challenges'' section, and it is correct.
The consequence: the exact packaged-build settings for the deliverable cannot be
recovered from the repository, only \textbf{inferred} from its two siblings. In
particular \code{console=True} is wrong for a GUI application (it would open a
black terminal window behind the tkinter window), so the real GUI spec almost
certainly used \code{console=False}. Marking that explicitly \textbf{inferred}:
nothing in the repo confirms it.

There is a second, sharper reason to care. The GUI is riddled with debug
\code{print} calls. Under \code{console=False} those prints go nowhere, and the
\code{error} widget is only ever written on success. So the packaged client build,
as best it can be reconstructed, had \textbf{no channel at all} for reporting a
failure to its user. That is consistent with the previous section's finding: a
matcher that silently produces zeros in a build that cannot report anything is a
tool that can be wrong for a long time without anyone learning.
\end{honest}

\subsection{Dependencies}

\begin{center}
\begin{tabular}{@{}lll@{}}
\toprule
\textbf{Package} & \textbf{Pin} & \textbf{Why} \\
\midrule
\code{pandas} & \code{>=2.0,<3.0} & Everything. Upper bound is load-bearing (below). \\
\code{openpyxl} & \code{>=3.1,<4.0} & The XLSX engine pandas calls for \code{.xlsx} files. \\
\code{tkinter} & (stdlib) & GUI and file pickers. Not pip-installable. \\
\bottomrule
\end{tabular}
\end{center}

\begin{keyidea}{Why \code{pandas<3} is not a lazy pin}
Chained assignment (\code{df[col][row] = v}) and writing into filtered views are
used throughout. Under pandas 3.0's copy-on-write these write to a temporary copy
and the result is discarded, \textbf{without an error}. The tools would run, report
success, and produce files with empty columns. The pin converts a silent
correctness failure into a dependency constraint. It is the right call for code
that is not going to be rewritten, and the README is honest that rewriting the
loops as a \code{merge} is the actual fix.
\end{keyidea}

Note that \code{openpyxl} is required by the GUI only; the two CLI tools are pure
CSV and never touch it.

\subsection{Security review}

\begin{keyidea}{Nothing to leak, and nothing leaked}
This project has no credentials by construction: no API client, no tokens, no shop
domain, no network calls of any kind. \file{.env.example} says so explicitly and
\file{.gitignore} excludes \file{.env} regardless. A scan of the working tree
found no keys, tokens, passwords or endpoints, and the git history was checked:
the client's real stock list and product export were \textbf{never committed}, so
there is nothing to scrub from history. The sample data is fabricated and the
README says so. This is a clean repository, and the offline design is why.
\end{keyidea}

\section{Honest overall assessment}

\subsection{What is genuinely good}

\begin{itemize}
  \item \textbf{The brand-filter insight.} Recognising that the join is ambiguous
  and that \code{Brand} is the disambiguator is the piece of thinking the project
  turns on. Filtering before joining removed a class of wrong match that no amount
  of fuzzy title matching would have fixed. Verified: it works.
  \item \textbf{The \code{Default Title} split.} Correctly reading Shopify's
  export conventions, including that variant rows have an empty \code{Title},
  is not obvious from the outside and is right.
  \item \textbf{The \code{total width} fallback and the series-2xxxx cage rule.}
  Both are fingerprints of someone who read the client's real data instead of
  assuming its shape.
  \item \textbf{Taking delivery seriously.} The GUI, the executables, the file
  pickers and the non-overwrite counter exist because the user was a person, not a
  terminal. The README is right that this took as much thought as the matching.
  \item \textbf{The sample data.} Rebuilding a fabricated fixture with the same
  column layout, the same HTML structure, and deliberate brand collisions, so the
  tools still run and can be checked without exposing client data, is better than
  most portfolio scrubs manage.
\end{itemize}

\subsection{What is weak}

\begin{itemize}
  \item \textbf{The GUI's quantity matcher does not match} on the shipped sample
  (verified). This is the headline finding.
  \item \textbf{Triplicated logic.} Three copies of two algorithms, already drifted
  in at least two observable ways (the dtype bug, the string-versus-float column).
  \item \textbf{No tests.} The sample data is a perfectly good fixture and there is
  not a single assertion against it. Any one test on the GUI's matcher would have
  caught the headline bug in 2023.
  \item \textbf{Debug \code{print}s throughout}, in the version whose likely
  packaging discards stdout.
  \item \textbf{Error handling.} \code{save\_data} catches only
  \code{FileNotFoundError}; \code{open\_file}'s retry loop cannot retry; nothing
  writes failures to the status widget.
  \item \textbf{Quadratic matching} and a full spreadsheet parse used as an
  existence check.
  \item \textbf{Brittle path handling}, splitting on a hard-coded \code{'/'}.
  \item \textbf{Brittle extraction.} The line-pair assumption is the store's
  template, not a guarantee, and the README says so.
\end{itemize}

\subsection{Verification status}

\begin{center}
\begin{tabular}{@{}ll@{}}
\toprule
\textbf{Claim} & \textbf{Status} \\
\midrule
CLI quantity/cost merge picks HI-TEC rows & \textbf{Verified by running it} \\
CLI merge matches variants via \code{Option1 Value} & \textbf{Verified by running it} \\
CLI dimensions extractor recovers all three & \textbf{Verified by running it} \\
Dimensions extractor drops description-less rows & \textbf{Verified by running it} \\
GUI quantity matcher matches nothing on the sample & \textbf{Verified by running its logic} \\
GUI window builds and matches the screenshot & \textbf{Verified} (code against \file{app.png}) \\
No credentials in tree or history & \textbf{Verified} \\
GUI's real \file{.spec} used \code{console=False} & \textbf{Inferred}, not in repo \\
Shopify import ignores absent rows & \textbf{Inferred}, not tested \\
GUI bug's effect on the real 2023 data & \textbf{Unresolvable} from this repo \\
1,600 stock rows / 11,000-line export & \textbf{Unverifiable}, data never committed \\
End-to-end live store round trip (2023) & \textbf{Not re-verified}, as README says \\
\bottomrule
\end{tabular}
\end{center}

\section{If it were rebuilt}

\begin{enumerate}
  \item \textbf{One module, two frontends.} A \file{core.py} holds the three
  operations as pure functions, each mapping a DataFrame to a DataFrame:
  \code{merge\_quantities} for tool 1, then \code{extract\_dimensions} for tool 2,
  then \code{detect\_variants} for the third. The CLI and the tkinter app both
  import it. The headline bug in this repository is a direct consequence of
  not doing this, which makes it the first thing to change.
  \item \textbf{Replace the loops with a \code{merge}.} One
  \code{pandas.merge} against the brand-filtered frame replaces both loops, is
  linear rather than quadratic, removes every chained assignment, and lifts the
  \code{pandas<3} pin as a side effect.
  \item \textbf{Normalise join keys explicitly, on both sides.} The GUI bug is a
  half-done version of exactly this: it normalised one side and not the other.
  Cast both to trimmed lowercase strings, or neither.
  \item \textbf{Distinguish ``no match'' from ``zero''.} A blank or \code{NaN},
  never a \code{0}, in a column that feeds a live storefront's inventory.
  \item \textbf{Test against the sample.} It is already an adversarial fixture.
  Six assertions would pin every behaviour this document had to verify by hand.
  \item \textbf{A dict, not parallel lists,} for the variant code tables.
  \item \textbf{\code{pathlib} for paths, \code{os.path.exists} for existence,
  \code{logging} instead of \code{print},} and a status widget that is written on
  failure as well as success.
\end{enumerate}

\begin{keyidea}{The fair summary}
Three small tools that solved a real and genuinely awkward problem for a real
client, under a constraint (no API) that forced text matching on data with no
shared key. The central idea is sound and verifiably works in the CLI tools. The
delivered GUI is a copy-paste fork of them that carries at least one serious,
silent bug, and the absence of both tests and a shared module is why that bug
could exist unnoticed. The code is 2023 freelance work whose weaknesses the author
has since named accurately in his own README; this document adds the one he had
not found, and it is the most interesting thing here.
\end{keyidea}

\end{document}
